\documentclass[book.tex]{subfiles}
\begin{document}

\chapter{Polynomial Interpolation}
\label{chap:polynomial}
\noindent\rule{11cm}{0.4pt} \\
\textbf{Prerequisites:} \autoref{chap:linear_systems} \\
\textbf{Difficulty Level:} * \\
\noindent\rule{11cm}{0.4pt}
\vspace{5mm}

In this chapter, we will introduce the two notions of curve fitting and interpolation and explain the difference between the two methods. We explain how solving polynomial interpolation with simultaneous equations is an ill-conditioned problem. We will then study how to perform polynomial interpolation using Newton's interpolating polynomial.

\section{Curve Fitting vs. Interpolation}

When data points are assumed to be approximate, we do not require that our resulting curve passes exactly through them, but the idea is to get as close as possible on average. In contrast, interpolation assumes the data points are known exactly or at least to very good precision, and a curve is chosen that passes through each of the data points instead of just as near as possible as in the least squares approach. If we have $n$ data points we wish to interpolate, either we can fit a polynomial of degree $n - 1$ or find $n - 1$ line segments which connect the dots or use splines.

\section{Generalized Least Squares}

In the last chapter, we have seen how to fit a linear model over the entire range of data. This resulted in the model predicting well in certain parts of the feature-space while performing poorly in other parts. Now, we'll try to do similar modeling, while dividing the state space into a finite number of buckets (or bins). This enables us to fit independent models for each bucket, thereby increasing accuracy.

\begin{example} The housing prices with additional features like Area, Number of Rooms, Lot Area, etc are given. Try to improve the accuracy by modeling three linear models as following

Divide the data into three bins as
\begin{enumerate}
\item For Price of house less than 5 percent
\item For price of house between 5 and 10 percent
\item For price of house greater than 10 percent
\end{enumerate}

Fit individual models for each bin, such that the classification accuracy is at-least 90 percent for each bin starting with using just one feature (or) until you run out of features.

\textit{Output: } 
\begin{Verbatim}
Without Different models for different bins: 
RMS Error: 0.034287
RMS Error in Dollars: 261440.790301
Mean Guess RMS Error in Dollars: 367118.703181
Ratio of RMS Error: 0.712142
Accuracy =
0.6418
0.5972
0.5843

With Different models for different bins: 
Number of features chosen for bin 1: 2
Number of features chosen for bin 2: 5
Number of features chosen for bin 3: 5
Accuracy =
0.9997
0.9384
0.9116
RMS Error: 0.023483
RMS Error in Dollars: 179061.680696
Mean Guess RMS Error in Dollars: 367118.703181
Ratio of RMS Error: 0.487749
\end{Verbatim}


%\textit{Matlab Code: } 
%\begin{Verbatim}
%
%load('Housing_Prices.mat')
%% Features Used
%% Feature 1 - Living Area
%% Feature 2 - Year Built
%% Feature 3 - Lot_Area
%% Feature 4 - Rooms
%% Feature 5 - Condition (A number between 1 to 5 assesing the condition of the house)
%
%%% Without normalization
%x1 = Area; x2 = Year_Built; x3 = Lot_Area; x4 = Rooms; x5 = Condition; y = Price;
%
%%% With normalization
%min_x1 = min(x1); max_x1 = max(x1);
%min_x2 = min(x2); max_x2 = max(x2);
%min_x3 = min(x3); max_x3 = max(x3);
%min_x4 = min(x4); max_x4 = max(x4);
%min_x5 = min(x5); max_x5 = max(x5);
%min_y = min(y); max_y = max(y);
%
%x1 = (x1 - min_x1)/(max_x1- min_x1);
%x2 = (x2 - min_x2)/(max_x2- min_x2);
%x3 = (x3 - min_x3)/(max_x3- min_x3);
%x4 = (x4 - min_x4)/(max_x4- min_x4);
%x5 = (x5 - min_x5)/(max_x5- min_x5);
%y = (y - min_y)/(max_y-min_y);
%
%X = [x1 x2 x3 x4 x5]; 
%
%%% Solving for the model
%num_features = 1;  % Starting with just 1 feature
%max_features = 5;
%[coeffs, Z] = linear_reg(X, y, num_features, false);
%coeffs
%y_actual = y;
%y_estimates = Z*coeffs;
%
%RMS_Error = sqrt(sum((y - y_estimates).^2)/num_data);
%fprintf("RMS Error: %f\n", RMS_Error);
%fprintf("RMS Error in Dollars: %f\n", RMS_Error*(max_y-min_y));
%
%RMS_Error2 = sqrt(sum((y - mean(y)).^2)/length(y));
%%fprintf("RMS Error: %f\n", RMS_Error2);
%fprintf("Mean Guess RMS Error in Dollars: %f\n", RMS_Error2*(max_y-min_y));
%ratio = RMS_Error/RMS_Error2;
%fprintf("Ratio of RMS Error: %f\n", ratio);
%
%figure
%plot(y_estimates,y_actual,'o');
%
%%% Separating the data into three bins
%% Bin 1 - Price less than 5 percent of maximum price
%% Bin 2 - Price between 5 and 10 percent of maximum price
%% Bin 3 - Price greater than 10 percent of maximum price
%fprintf("\nWith Different models for different bins: \n\n");
%
%y1_ind = y >= 0 & y <= 0.05;
%y2_ind = y > 0.05 & y <= 0.1;
%y3_ind = y > 0.1;
%
%y1 = y_actual(y1_ind);
%y2 = y_actual(y2_ind);
%y3 = y_actual(y3_ind);
%X1 = [x1(y1_ind) x2(y1_ind) x3(y1_ind) x4(y1_ind) x5(y1_ind)]; 
%X2 = [x1(y2_ind) x2(y2_ind) x3(y2_ind) x4(y2_ind) x5(y2_ind)]; 
%X3 = [x1(y3_ind) x2(y3_ind) x3(y3_ind) x4(y3_ind) x5(y3_ind)]; 
%
%Accuracy = zeros(3,1);
%y1_estimates = y_estimates(y1_ind);
%y2_estimates = y_estimates(y2_ind);
%y3_estimates = y_estimates(y3_ind);
%y1_est_ind = y1_estimates >= 0 & y1_estimates <= 0.05;
%Accuracy(1) = nnz(y1_est_ind)/length(y1_estimates);
%y2_est_ind = y2_estimates > 0.05 & y2_estimates <= 0.1;
%Accuracy(2) = nnz(y2_est_ind)/length(y2_estimates);
%y3_est_ind = y3_estimates > 0.1;
%Accuracy(3) = nnz(y3_est_ind)/length(y3_estimates);
%Accuracy
%
%%% First bin
%cur_accuracy = Accuracy(1);
%cur_features = num_features;
%while (cur_accuracy < 0.9 && cur_features < max_features)
%cur_features = cur_features + 1;
%[coeffs, Z1] = linear_reg(X1, y1, num_features, false);
%y_temp_estimates = Z1*coeffs;
%y1_ind = y1 >= 0 & y1 <= 0.05;
%y1_estimates = y_temp_estimates(y1_ind);
%y1_est_ind = y1_estimates >= 0 & y1_estimates <= 0.05;
%cur_accuracy = nnz(y1_est_ind)/length(y1_estimates);
%Accuracy(1) = cur_accuracy;
%end
%
%fprintf("Number of features chosen for bin 1: %d\n", cur_features);
%%% Second bin
%cur_accuracy = Accuracy(2);
%cur_features = num_features;
%% y2 = y2 - Price_bin1;
%
%while (cur_accuracy < 0.9 && cur_features < max_features)
%cur_features = cur_features + 1;
%[coeffs, Z2] = linear_reg(X2, y2, num_features, true);
%y_temp_estimates = Z2*coeffs;
%y2_ind = y2 > 0.05 & y2 <= 0.1;
%y2_estimates = y_temp_estimates(y2_ind);
%y2_est_ind = y2_estimates > 0.05 & y2_estimates <= 0.1;
%cur_accuracy = nnz(y2_est_ind)/length(y2_estimates);
%Accuracy(2) = cur_accuracy;
%end
%
%fprintf("Number of features chosen for bin 2: %d\n", cur_features);
%%% Third bin
%cur_accuracy = Accuracy(3);
%cur_features = num_features;
%
%while (cur_accuracy < 0.9 && cur_features < max_features)
%cur_features = cur_features + 1;
%[coeffs, Z3] = linear_reg(X3, y3, num_features, true);
%y_temp_estimates = Z3*coeffs;
%y3_ind = y3 > 0.1;
%y3_estimates = y_temp_estimates(y3_ind);
%y3_est_ind = y3_estimates > 0.1 ;
%cur_accuracy = nnz(y3_est_ind)/length(y3_estimates);
%Accuracy(3) = cur_accuracy;
%end
%fprintf("Number of features chosen for bin 3: %d\n", cur_features);
%\end{Verbatim}
\end{example}

\section{Introduction to Interpolation}
Polynomial interpolation is the most common method used to estimate the intermediate values between precise data points. The general formula for an $(n-1)$th order polynomial can be written as 
\begin{align}
f(x) = p_1x^{n-1} + p_2x^{n-2} +  . . . + p_{n-1}x + p_n
\vspace{-8pt}
\end{align}
For $n$ data points, there is one and only one polynomial of order $(n-1)$ that passes through all the points. Polynomial interpolation consists of determining the unique $(n-1)$th order polynomial that fits $n$ data points. This polynomial then provides a formula to compute intermediate values.
\begin{marginfigure}
	\centering
	\includegraphics[width = 2.4in]{figures/part1b/data_interpolation/interpolation.PNG}
	\centering
	\caption{Examples of interpolating polynomials:(a) first-order (linear)
		(b) second-order (quadratic or parabolic) (c) third-order (cubic)}
	\label{fig:1}
\end{marginfigure}
\subsection{Determining Polynomial Coefficients}
We need $n$ data points to determine the $n$ coefficients. This allows us to generate $n$ linear algebraic equations that we can solve simultaneously for the
coefficients.\\
\begin{example}
Determine the coefficients of the parabola $f(x)$ which passes through the points: (-2,-1),(-0.5,-10) and (1,62).
\begin{align}
f(x) = p_1x^2 + p_2x + p_3
\vspace{-8pt}
\end{align}
Each of these pairs can be substituted into above equation to yield a system of three equations:

\begin{align}
\begin{gathered}
-1 = p_1(-2)^2 + p_2(-2) + p_3, -10 = p_1(-0.5)^2 + p_2(-0.5) + p_3\\
62 = p_1(1)^2 + p_2(1)x + p_3
\end{gathered}
\end{align}
or in matrix form as shown on the right.
\begin{marginfigure}
\begin{align}
\begin{bmatrix}
4 & -2 & 1 \\ 
0.25 & -0.5 & 1 \\ 
1 & 1 & 1 
\end{bmatrix}
\left \{
\begin{tabular}{c}
$p_1$ \\
$p_2$ \\
$p_3$ 
\end{tabular}
\right \}
=
\left \{
\begin{tabular}{c}
-1 \\
-10 \\
62 
\end{tabular}
\right \}
\end{align}
\end{marginfigure}
Thus, the problem reduces to solving three simultaneous linear algebraic equations for the three unknown coefficients. When solved in MATLAB we get $p_1 = 18$, $p_2 = 39$ and $p_3 = 5$. Thus, the parabola that passes exactly through the three points is
\begin{align}
f(x) = 18x^2 + 39x + 5
\end{align}

\end{example}
Although the approach in the above example provides an easy way to perform interpolation, it has a serious deficiency. To understand this flaw, notice that the coefficient matrix has a decided structure which can be seen clearly by expressing it in general terms. 
\begin{marginfigure}
\[
\begin{bmatrix}
x_1^2 & x_1 & 1 \\ 
x_2^2 & x_2 & 1 \\  
x_3^2 & x_3 & 1 
\end{bmatrix}
\left \{
\begin{tabular}{c}
$p_1$ \\
$p_2$ \\
$p_3$ 
\end{tabular}
\right \}
=
\left \{
\begin{tabular}{c}
$f(x_1)$ \\
$f(x_2)$ \\
$f(x_3)$ 
\end{tabular}
\right \}
\vspace{-6pt}
\]
\end{marginfigure}
Coefficient matrices of this form are referred to as $\textit{Vandermonde matrices}$. Such matrices are very ill-conditioned. That is, their solutions are very sensitive to round-off errors. The ill-conditioning becomes even worse as the number of simultaneous equations becomes larger. In this lecture, we will describe an alternative that is well-suited for
computer implementation, the Newton polynomials.\\
%\subsection{MATLAB Functions: $\textrm{polyfit}$ and $\textrm{polyval}$}
%Recall that the $\verb|polyfit|$ function can be used to perform polynomial regression. In such applications, the number of data points is greater than the number of coefficients being estimated. Consequently, the least-squares fit line does not necessarily pass through any of the points, but rather follows the general trend of the data. For the case where the number of data points equals the number of coefficients, $\verb|polyfit|$ performs interpolation. That is, it returns the coefficients of the polynomial that pass directly through the data points. For example, it can be used to determine the coefficients of the parabola that passes through the three points from the above example:
%\begin{verbatim}
%>> x = [-2 -0.5 1]; y = [1 -10 62];
%>> p = polyfit(x,y,2)
%p = 18.0000   39.0000    5.0000
%\end{verbatim}
%We can then use the $\verb|polyval|$ function to perform an interpolation (i.e. find the value of the function at any intermediate point).
%\begin{verbatim}
%>> d = polyval(p,20)
%d = 7.9850e+03
%\end{verbatim}

\section{Newton Interpolating Polynomial}

\subsection{Linear Interpolation}
The simplest form of interpolation is to connect two data points with a straight line, depicted graphically in Fig.\ref{fig:2}. Using similar triangles,

\begin{marginfigure}[-5\baselineskip] 
	\centering
	\includegraphics[width = 2.4in]{figures/part1b/data_interpolation/linear.png}
	\centering
	\caption{Graphical depiction of linear interpolation.}
	\label{fig:2}
\end{marginfigure}
\begin{align}
\dfrac{f_1(x)-f(x_1)}{x-x_1} = \dfrac{f(x_2)-f(x_1)}{x_2-x_1}
\end{align}
which can be rearranged to yield

\begin{align}
f_1(x) = f(x_1) + \dfrac{f(x_2)-f(x_1)}{x_2-x_1}(x-x_1) 
\end{align}
which is the \textit{Newton linear-interpolation} formula. The notation $f_1(x)$ designates that this is a first-order interpolating polynomial. Notice that besides representing the slope of the line connecting the points, the term 
\begin{align}
[f(x_2) - f(x_1)]/(x_2 - x_1)
\end{align} 
is a finite-difference approximation of the first derivative. In general, the smaller the interval
between the data points, the better the approximation. As the interval decreases, a continuous function will be better approximated by a straight line.
\subsection{Quadratic Interpolation}
The error in the example resulted from approximating a curve with a straight line. Consequently, a strategy for improving the estimate is to introduce some curvature into the line
connecting the points. If three data points are available, this can be accomplished with a second-order polynomial (also called a quadratic polynomial or a parabola). A particularly convenient form for this purpose is
\begin{align}
f_2(x) = b_1 + b_2(x-x_1) + b_3(x-x_1)(x-x_2)
\end{align}
with $x = x_1$, $b_1$ could be evaluated as $b_1 = f(x_1)$. This is substituted in the main equation which is evaluated at $x = x_2$, to get
\begin{align}
b_2 = \dfrac{f(x_2)-f(x_1)}{x_2-x_1}
\end{align}
Finally $b_1$ and $b_2$ can be substituted into the main equation which is evaluated at $x = x_3$ and solved for
\begin{align}
b_3 = \dfrac{\dfrac{f(x_3)-f(x_2)}{x_3 - x_2} - \dfrac{f(x_2)-f(x_1)}{x_2 - x_1}}{x_3-x_1}
\end{align}
Notice that, as was the case with linear interpolation, $b_2$ still represents the slope of the line connecting points $x_1$ and $x_2$. Thus, the first two terms of the main equation are equivalent to linear interpolation between $x_1$ and $x_2$. The last term, $b_3(x-x_1)(x-x_2)$, introduces the second-order curvature into the formula. Thus the terms are added sequentially to capture increasingly higher-order curvature.
\subsection{General form of Newton's Interpolating Polynomials}
The preceding analysis can be generalized to fit an $(n-1)$th-order polynomial to $n$ data points. The $(n-1)$th-order polynomial is
\begin{align}
f_{n-1}(x) = b_1 + b_2(x-x_1) + ... + b_n(x-x_1)(x-x_2)...(x-x_{n-1})
\end{align}
As was done previously with linear and quadratic interpolation, data points can be used to evaluate the coefficients $b_1, b_2, . . . , b_n$ . For an $(n - 1)$th-order polynomial, $n$ data points are required.
\begin{align}
b_1 &= f(x_1)\\
b_2 &= f[x_2,x_1] \\
b_3 &= f[x_3,x_2,x_1], ... \\
b_n &= f[x_n,x_{n-1},...,x_2,x_1]
\end{align}
where the bracketed function evaluations are finite divided differences.
\begin{marginfigure}
	\centering
	\includegraphics[width = 2.4in]{figures/part1b/data_interpolation/polynomial.png}
	\centering
	\caption{Graphical depiction of the recursive nature of finite divided differences.}
	\label{fig:4}
\end{marginfigure}
 For example, the $n$th finite divided difference is
\begin{align}
f[x_n,x_{n-1},...,x_2,x_1] = \dfrac{f[x_n,x_{n-1},...,x_2] - f[x_{n-1},x_{n-2},...,x_1]}{x_n - x_1} 
\end{align}
These differences can be used to evaluate the coefficients,
which can then be substituted into the main equation to yield the general form of Newton's interpolating polynomial:
\vspace{-8pt}
\begin{align}
\begin{gathered}
f_{n-1}(x) = f(x_1) + (x - x_1)f[x_2,x_1] + (x - x_1)(x -  x_2)f[x_3,x_2,x_1]+\\ ... + (x - x_1)(x - x_2)...(x - x_{n-1})f[x_n,x_{n-1},...,x_2,x_1] 
\end{gathered}
\end{align}
We should note that it is not necessary that the data points used be equally spaced or that the abscissa values necessarily be in ascending order. \\
\begin{example} You are an engineer at Tesla and you want to estimate the Coefficient of drag ($C_d$) and Coefficient of rolling resistance ($C_{rr}$) given the power v/s velocity profile of a real car. The load power on the vehicle can be calculated based on the velocity profile v(t) as
\begin{equation*}
P(t) = F_{total}.v(t) = \frac{1}{2}\rho C_dA.v(t)^3 + C_{rr}mg.v(t) + \alpha.v(t)^2 + \beta
\end{equation*}
where $P(t)$ is the load power, $\rho$ is the density of air, $A$ is the frontal area of the vehicle, $v(t)$ is the instantaneous velocity, $m$ is the total mass of the vehicle, $\alpha$ and $\beta$ are constants and $g$ is acceleration due to gravity. Given: $\rho = 1.17 kg/m^3$, $A = 5.5 m^2$, $m = 1600 kg$, $g = 9.81 m/s^2$ and the values of power and velocity at four points as (0,0), (12,3858.98),(28.667,36725.26), (35.334,66608.53).  Evaluate the power at $v_t = 22 m/s$, $C_d$ and $C_{rr}$ using Newton interpolation. 

\textit{Solution}: \textit{MATLAB implementation of Newton Interpolation}: 
\begin{marginfigure}[5\baselineskip]
\begin{Verbatim}
v = [0; 12; 28.667; 35.334]; 
p = [0; 3858.98; 36725.26; 66608.53];
v_int = 22; 
\end{Verbatim}
\end{marginfigure}
\begin{Verbatim}
mass = 1600; rho = 1.17; A_F = 5.5; g = 9.81; v_int = 22;
v = [0; 12; 28.667; 35.334]; %m/s
p = [0; 3858.98; 36725.26; 66608.53]; %Watt
n = length(v); b = zeros(n,n);b(:,1) = p(:); 
for j = 2:n
	for i = 1:n-j+1
		b(i,j) = (b(i+1,j-1)-b(i,j-1))/(v(i+j-1)-v(i))
vt = 1; p_int = b(1,1);
for j = 1:n-1
	vt = vt*(v_int-v(j));
	p_int = p_int + b(1,j+1)*vt;
end
C_D = (b(1,4))/(0.5*rho*A_F); # Coefficient of Drag
# Coefficient of Rolling Resistance
C_R = (b(1,2) - b(1,3)*(v(2) + v(1)) ...
+ b(1,4)*(v(1)*v(2) + v(2)*v(3) + v(1)*v(3)))/(mass*g);
\end{Verbatim}
\begin{marginfigure}[-5\baselineskip]
	\textit{Output}: 
	\begin{Verbatim}
    p_int = 1.7664e+04
    C_D = 0.4399
    C_R = 0.0075
	\end{Verbatim}
\end{marginfigure}
Solving the $(n-1)$th-order polynomial (here $n=4$) for the coefficient of $v_t^3$ and $v_t$ and comparing with the equation for Power $P(t)$ we get,
\begin{align}
\begin{gathered}
b_4 = \frac{1}{2}\rho C_dA.v(t)^3 \\
b_2 - b_3(v_2 + v_1) + b_4((v_1 \times v_2) + (v_2 \times v_3) + (v_1 \times v_3)) = C_{rr}mg.v(t) 
\end{gathered}
\end{align}
%\end{document}
\end{example}
\section{Exercises}
\begin{enumerate}
    \item TODO
\end{enumerate}

\end{document}